\textbf{IV -- GLI AVAMPOSTI}

Per avamposto si intende il pezzo (pedone compreso) più avanzato dello
schieramento.

Per intenderci meglio, già la prima mossa, se si tratta di una spinta di
pedone di due passi, secondo la concezione di Canal (1896--1981), cui
qui ci rifacciamo, introduce sulla scacchiera un avamposto. Quando è un
pezzo non facilmente scacciabile o è un pedone, l'avamposto diviene il
leitmotiv della partita e su di esso si concentrano i pezzi dell'uno e
dell'altro schieramento.

\textbf{a) Il caso del pedone}

Nel capitolo precedente abbiamo già avuto modo di constatare l'efficacia
del pedone come avamposto; la sua importanza dipende dalla mobilità e
dal controllo esercitato sulle case limitrofe. Per quanto riguarda
quest'ultimo aspetto vale la pena soffermarsi un momento sul diagramma
che segue:

\bookdiagramplaceholder{assets/diagrams/image70.png}{70}

-- la casa e6 è la casa di spinta. Chi si difende deve averne il
controllo per evitare pericolose spinte di pedone. Nella casa di spinta
un giocatore tenderà a piazzare un pezzo bloccatore mentre l'altro si
adopererà a rimuoverlo.

-- le case d6 e f6 sono case di controllo. Se ne ricorrono le
condizioni, possono esservi installati pezzi, in particolare Cavalli.

-- le case d5 e f5 sono case di affiancamento. Affiancare un pedone con
un secondo pedone ne aumenta la mobilità. Queste case sono spesso
sfruttate dagli Alfieri amici, che non trovano intralcio alla loro
azione dal pedone e contribuiscono al controllo della casa di spinta.

-- infine d4, e4 e f4 sono case di sostegno. Quelle laterali, se
occupate da pedoni, sostengono l'avamposto e minacciano di affiancare un
secondo pedone al primo, se occupate da Cavalli controllano l'importante
casa di spinta; quella centrale (ma più che di casa andrebbe parlato di
colonna) è occupata di preferenza dalla Torre che sostiene, specie nel
finale, l'avanzata del pedone.

\textbf{b) Il caso del pezzo.}

Un pezzo può essere l'avamposto dell'armata solo se piazzato su casa
debole o non controllabile senza causare danni di altro tipo alla
propria posizione.

Si consideri questo schema:

\bookdiagramplaceholder{assets/diagrams/image71.png}{71}

La struttura di pedoni del Nero non presenta debolezze. In queste
condizioni non dovrebbe essere possibile al Bianco creare un pezzo come
avamposto oltre la quarta traversa. Invece egli può giocare ¤c3-d5 senza
temere di essere scacciato perché la spinta del pedone nero da c7 in c6
indebolirebbe il §d6.

Una bella partita dove un avamposto nero sfrutta una casa debole nel
cuore dello schieramento avversario è la seguente:

Karpov--Kasparov, Mosca, 16\textsuperscript{a} del match mondiale 1985
\textbf{(Difesa Siciliana)}

\textbf{l.e4 c5 2.¤f3 e6 3.d4 ¤xd4 4.¤xd4 ¤c6 5.¤b5 d6 6.c4} negli anni
Settanta Fischer preferiva giocare 6.¥f4 e5 7.¥e3 ¤f6 8.¥g5 con rientro
nella variante Svesnhikov della Siciliana, ma con un tempo in meno per
il Bianco. Egli considerava tanto gravi le debolezze del Nero (buco in
d5 e pedone debole in d6) da non preoccuparsi troppo della pe¢dita di
tempo. Il sistema scelto da Karpov, noto col nome di Maroczy, è molto
solido e tende al controllo della casa d5.

\textbf{6... ¤f6 7.¤1c3 a6 8.¤a3 d5!?} Kasparov aveva già tentato questa
sorprendente spinta nella dodicesima partita del match; questo
sacrificio di pedone era stato considerato dubbio dai teorici e pochi
pensavano a una replica.

\textbf{9.¤xd5 exd5 10.exd5} il Bianco ritiene che 10.¤xd5 ¤xd5 11.£xd5
¥xe6 12.£xd8+ ¢xd8 concede troppo gioco all'avversario.

\textbf{10... ¤b4 11.¥e2} nell'occasione precedente Karpov aveva giocato
11.¥e4.

\textbf{11... ¥c5} sembra che questa mossa sia sfuggita al team di
Karpov che, forse, si aspettava 11... ¤bxd5 12.¤xd5 ¤xd5 13.0-0 ¥e7
14.¥f3 ¥e6 con un leggero vantaggio del Bianco dopo la manovra ¤c2-d4.

\textbf{12.0-0 0-0 13.¥f3} si tiene stretto il pedone di vantaggio che
invece era meglio restituire con 13.¥g5 ¤bxd5 14.¤xd5 £xd5 15.¥xf6 £xd1
16.¦fxd1 gxf6 con pari possibilità.

\textbf{13... ¥f5 14.¥g5 ¦e8} 14... ¥d3 è confutata da 15.¦e1 ¤g4
16.¥xd8 ¤xf2 17.£d2!

\textbf{15.£d2} non è possibile giocare 15.¦e1 per 15... ¦xe1+ 16.£xe1
¤d3

\textbf{15... b5 16.¦ad1 ¤d3!}

\bookdiagramplaceholder{assets/diagrams/image72.png}{72}

Il commento più colorito a questa mossa lo hanno scritto Keene e Goodman
in ``Karpov--Kasparov: atto secondo'', edito in Italia dalla Mursia nel
1985: \emph{``Questo pezzo comincia come Cavallo ma si trasforma ben
presto in una piovra, posta al centro della scacchiera, i cui tentacoli
si estendono in ogni direzione e si protendono sulle caselle chiave
della posizione del Bianco''.} Nella casa d3 il Nero ha piazzato un
avamposto che condiziona in modo pesante l'intera posizione.

\textbf{17.¤ab1} era minacciata la forchetta in b4. A 17.¤c2 poteva
seguire 17... ¤xb2 18.¤e3 ¥xe3 19.fxe3 ¤xd1 20.¥xd1 £d7.

\textbf{17... h6 18.¥h4 b4 19.¤a4 ¥d6 20.¥g3 ¦c8 21.b3} per far
rientrare in gioco il ¤a4 tramite la casa b2.

\textbf{21... g5!} schioda il ¤f6 per portarlo a sostegno del ¤d3.

\textbf{22.¥xd6 £xd6 23.g3} non va 23.¤b2 ¤xb2 24.£xb2 g4 25.¥e2 ¦c2 e
vince.

\textbf{23... ¤d7!} libera la casa f6 alla Donna e prepara un altro
Cavallo per la casa d3 nel caso il primo venisse cambiato.

\textbf{24.¥g2} il ¤a4 non può ancora essere riportato in gioco tramite
b2, seguirebbe: 24... ¤7e5 25.¥g2 ¤xb2 26.£xb2 ¦c2 27.£d4 ¦xa2 con
vantaggio evidente.

\textbf{24... £f6 25.a3 a5 26.axb4 axb4 27.£a2 ¥g6} vale la pena notare
come l'avamposto in d3 riduca al minimo le possibilità del Bianco.

Secondo le analisi di Suetin, pubblicate dalla ``Pravda'', era sbagliata
27... ¤f4 28.gxf4 ¦c2 29.¤b2 ¦xb2 30.fxg5 hxg5 31.£a4.

\textbf{28.d6 g4} il Bianco respira dopo 28... £xd6 29.¤d2 ¦c2 30.¤c4.

\textbf{29.£d2 ¢g7 30.f3} aprire le linee affretta la fine ma non si
vede come il Bianco possa resistere a manovre stritolanti del tipo di
h5-h4 e ¦h8.

\textbf{30... £xd6 31.fxg4 £d4+ 32.¢h1 ¤f6 33.¦f4 ¤e4} tutti i pezzi del
Nero partecipano all'attacco. Il resto è affidato alla tattica.

\textbf{34.£xd3} dà la Donna per tre pezzi leggeri, ma sono indigesti.

\textbf{34... ¤f2+ 35.¦xf2 ¥xd3 36.¦fd2 £e3 37.¦xd3 ¦c1 38.¤b2} se
38.¦xe3 ¦xd1+ e ¦xe3.

\textbf{38... £f2}

\bookdiagramplaceholder{assets/diagrams/image73.png}{73}

\textbf{39.¤d2 ¦xd1+ 40.¤xd1 ¦e1+ 41. il Bianco abbandona}. A 41.¤f1
¦xf1+ 42.¥xf1 £xf1 matto. Forse la sconfitta più cocente di tutta la
carriera scacchistica di Karpov.
